
\SGAStatement{Corollary 1.3.3}{%
If \(G\) is a smooth affine connected \(k\)-group, then the union of
\(T\cdot\operatorname{rad}^{u}(G)\), where \(T\) ranges over the maximal tori
of \(G\), contains a dense open subset of \(G\).}
\SGARefTarget{sga3:xix:corollary:1:3:3}

\par\medskip\noindent
\textbf{\textit{\SGARefLink{sga3:xix:notation:1:4:0}{Notation 1.4.0}.}}---%
\footnote{Editor's note: this notation, which appeared in \SGARefLink{sga3:xix:lemma:2-3}{2.3} of the original,
has been inserted here; moreover, the statement and proof of \SGARefLink{sga3:xix:lemma:1-4}{1.4} have been
expanded.}
\textit{We shall systematically use the following notation.  If \(S\) is a
scheme and \(s\) is a point of \(S\), then \(\bar{s}\) denotes the spectrum of
an algebraic closure \(\overline{\kappa(s)}\) of \(\kappa(s)\).}
\SGARefTarget{sga3:xix:notation:1:4:0}

\par\medskip\noindent
\SGAInlineRefTarget{sga3:xix:lemma:1-4}\textbf{\textit{\SGARefLink{sga3:xix:lemma:1-4}{Lemma 1.4}.}}---\footnotemark[5]
\textit{Let \(S\) be a scheme and let \(Q\) be an \(S\)-torus acting on an
\(S\)-group scheme \(H\), separated and smooth over \(S\).
\begin{enumerate}
\item[\SGAInlineRefTarget{sga3:xix:notation:1:4:0:item:i}(i)] The fixed-point functor \(H^{Q}\) is representable by a closed
  subscheme of \(H\), smooth over \(S\).
\item[(ii)] If, moreover, \(H\) is affine over \(S\) and has connected fibers,
  then the same holds for \(H^{Q}\).
\end{enumerate}}
\SGARefTarget{sga3:xix:lemma:1:4}

\SGAProof{Proof}
First, by \SGARefLink{sga3:viii:examples:6:5:item:d:01}{VIII 6.5(d)},%
\footnote{Editor's note: see also the additions made in
VI\textsubscript{B}, \SGARefLink{ex:sga3-VIB-6.2.4-d}{6.2.4(d)} and \SGARefLink{concl:sga3-VIB-6.5.5-a}{6.5.5(a)}.}
since \(H\) is separated over \(S\) and \(Q\) is essentially free over \(S\),
\(H^{Q}\) is representable by a closed subscheme of \(H\).  In particular, if
\(H\) is affine over \(S\), then so is \(H^{Q}\).

Now consider the semidirect product \(G=H\cdot Q\); it is smooth and separated
over \(S\).  The centralizer \(\underline{\operatorname{Centr}}_{G}(Q)\) is
then representable by a closed subscheme of \(G\), of finite presentation over
\(G\), by \SGARefLink{sga3:xi:corollary:6:11:item:a}{Exposé XI 6.11(a)},\footnotemark[6] and it is smooth over \(S\) by
\SGARefLink{sga3:xi:corollary:2:4}{Exposé XI 2.4}.

Since
\[
  \underline{\operatorname{Centr}}_{G}(Q)=H^{Q}\cdot Q,
\]
it follows that \(H^{Q}\) is smooth over \(S\).  Indeed, using the section
\[
  \sigma:H^{Q}\longrightarrow
  \underline{\operatorname{Centr}}_{G}(Q)
\]
deduced from the identity section \(\epsilon_{Q}\) of \(Q\), one sees at once
that \(H^{Q}\) is formally smooth over \(S\); and since \(\epsilon_{Q}\) is of
finite presentation, \(H^{Q}\) is locally of finite presentation over \(S\).
This proves \SGARefLink{sga3:xix:notation:1:4:0:item:i}{(i)}.

Finally, suppose that \(H\) is affine over \(S\) and has connected fibers.
Then every geometric fiber \(G_{\bar{s}}\) of \(G\) is a smooth affine
connected \(\kappa(\bar{s})\)-group.  Hence, by the first assertion of \SGARefLink{sga3:xix:paragraph:1-3}{1.3},
the group
\[
  \underline{\operatorname{Centr}}_{G_{\bar{s}}}(Q_{\bar{s}})
  =
  \bigl(\underline{\operatorname{Centr}}_{G}(Q)\bigr)_{\bar{s}}
  =
  (H^{Q})_{\bar{s}}\cdot Q_{\bar{s}}
\]
is also connected, and this implies that \((H^{Q})_{\bar{s}}\) is connected.
\hfill\(\square\)

\par\medskip\noindent\SGAInlineRefTarget{sga3:xix:paragraph:1-5}\textbf{1.5.}\quad
Recall that the \emph{reductive rank} of a smooth affine \(k\)-group \(G\) is
the common dimension of its maximal tori.  We denote it by
\(\operatorname{rgred}(G/k)\) or \(\operatorname{rgred}(G)\).  The equality
\(\operatorname{rgred}(G/k)=0\) holds if and only if \(G\) is unipotent
(that is, \(G=\operatorname{rad}^{u}(G)\)), by the \emph{Bible}, §6.4,
Corollary 1 to \SGARefLink{sga3:xix:lemma:1:4}{Theorem 4}, or [Ch05], §6.5, Corollary 1 to \SGARefLink{sga3:xix:theorem:2:5}{Theorem 5}.

If \(H\) is a normal subgroup of the smooth affine \(k\)-group \(G\), then the
quotient \(G/H\) is affine and smooth (Exposé VI\textsubscript{B}, \SGARefLink{thm:sga3-VIB-11.17}{11.17} and
VI\textsubscript{A}, \SGARefLink{thm:sga3-VIA-3.3.2:item:iii}{3.3.2(iii)}).  Moreover (\emph{Bible}, §7.3,
Theorem 3(a),(c)),
\[
  \operatorname{rgred}(G)
  =
  \operatorname{rgred}(G/H)+\operatorname{rgred}(H).
\]
